\section{The positive-control formula for the original moment limit}
\label{S:intrinsic-second-moment}

Retain every original covariance and amplitude.
Write $a=\alpha_0$, $I=(0,1)$, $K(s,t)=|s-t|^{-a}$ and
$H_I=H_t\otimes H_{\mathrm{sp}}$ for the ORIGINAL raw covariance.
Keep the uniform $E$ mean $\mathcal T_R=T_R(C_t\otimes I)$ and
$C_{\rm mean}$ of~\eqref{E:uniform-completed-mean}, with fixed
$2<r<2/a$. To distinguish it from the map below, write $J_R=T_R$.
The physical mask is $m_R=\eta(\cdot/R)^2$, used once.
Let $\mathcal C_{\mathrm{pos}}$ be the full nonseparable class
in~\eqref{E:variational-positive-class}.
For $\Psi$ in that class, retain its unmasked mean $h_{\Psi,s}$,
raw variance $v_\Psi$, integrated full-gradient rate $I_\Psi$ and
positive scores $S_c^{\mathrm{pos}}$ from
Section~\ref{S:unmasked-variational}.
No first-moment or author-constant interpretation of the label one
is assumed.

\begin{lemma}[Raw inclusion and the actual smooth mean]
  \label{L:positive-raw-embedding}
  For any real Hilbert space $G$ there is a bounded map
  $J_G:L^2(I;G)\to H_t\otimes G$ with
  \begin{equation}\label{E:positive-bochner-inclusion}
    \|J_GF\|^2
    =\int_{I^2}K(s,t)\langle F(s),F(t)\rangle_G\ud s\ud t
    \le K_0\|F\|_{L^2(I;G)}^2,\qquad K_0=\frac{2^a}{1-a}.
  \end{equation}
  The natural inclusion puts every $\Psi\in\mathcal C_{\mathrm{pos}}$
  in $H_I$, with
  \begin{equation}\label{E:positive-raw-mean}
    \|\Psi\|_{H_I}^2=v_\Psi,\qquad
    \mathcal T_R\Psi(s)=B_E(m_Rh_{\Psi,s})\quad\hbox{a.e. }s.
  \end{equation}
  For each fixed smooth $\Psi$ the right side gives a canonical
  version for every $R\ge1$ and every $s\in I$.
\end{lemma}

\begin{proof}
  Direct integration and concavity of $x^{1-a}$ give
  \[
    \int_I K(s,t)\ud t
    =\frac{s^{1-a}+(1-s)^{1-a}}{1-a}\le K_0 .
  \]
  For scalar $f\ge0$, Cauchy--Schwarz against $K(s,t)\ud t$ and then
  Tonelli give $\|Kf\|_2\le K_0\|f\|_2$.
  The same estimate for absolute values bounds the bilinear integral
  by $K_0\|f\|_2\|g\|_2$.
  Approximate scalar $L^2(I)$ functions by real $C_c^\infty(I)$
  functions. This bound makes their raw test-null classes Cauchy
  in $H_t$, defines the bounded inclusion, and identifies its pairing
  with the integral. Finite $G$-valued tensors then give
  \eqref{E:positive-bochner-inclusion}. Simple fields are dense in
  Bochner $L^2(I;G)$; the same bound on differences extends $J_G$.
  Absolute bilinear continuity extends the integral identity as well.

  For $G=H_{\mathrm{sp}}$, the Bochner integral
  $\int_IK(s,t)F(t)\ud t$ exists almost everywhere: its norm is
  bounded by $K\|F\|(s)$, which belongs to $L^2$.
  Simple-field approximation and the Schur estimate give a strongly
  measurable $L^2$ version and a bounded map of norm at most $K_0$.
  On smooth finite tensors it agrees with
  $(C_t\otimes I)J_G$ by the actual covariance definition.
  Approximation by such tensors converges in $L^2(I;G)$.
  The integral converges there by Schur; the completed tensor map
  converges in $L^r(I;G)$ by Lemma~\ref{L:uniform-mean-tensor}
  and~\eqref{E:positive-bochner-inclusion}, hence also in $L^2$.
  Uniqueness of the $L^2$ limit identifies these versions.
  In particular the integral belongs to $L^r$.

  Choose a closed ball $B$ containing all spatial supports of $\Psi$.
  For real smooth $q$ supported there, the physical covariance gives
  \begin{equation}\label{E:positive-local-spatial-bound}
    \|q\|_{H_{\mathrm{sp}}}^2
    \le |B|\,\Gamma(B-B)\|q\|_\infty^2.
  \end{equation}
  Indeed in $\int\ud x\int\Gamma(\ud y)q(x)q(x-y)$ only
  $x\in B,y\in B-B$ contribute. Absolute values and local finiteness
  prove the bound, including white and singular Radon covariances.
  Thus $t\mapsto\Psi_t$ is continuous, bounded and compactly supported
  as an $H_{\mathrm{sp}}$-valued field, and belongs to $L^2$.
  Formula~\eqref{E:positive-bochner-inclusion} and absolute Fubini give
  \[
    \|\Psi\|_{H_I}^2
    =\int_{I^2}K(s,t)\langle\Psi_s,\Psi_t\rangle_{H_{\mathrm{sp}}}
    \ud s\ud t
    =\int_{I\times\R^d}\Psi(s,x)h_{\Psi,s}(x)\ud x\ud s=v_\Psi.
  \]
  The compact spatial supports and~\eqref{E:positive-local-spatial-bound}
  make all these absolute integrals finite.

  Test the spatial row~\eqref{E:uniform-mean-row} first with displacement
  $f$ and dual second component $\Lambda g$, with $f,g$ compact smooth.
  The row product is $m_Rfg$, so the preceding temporal identity gives
  \[
    \left\langle\int_IK(s,t)\Psi_t\ud t,m_Rfg\right\rangle_{H_{\mathrm{sp}}}
    =\int m_Rh_{\Psi,s}fg .
  \]
  This Fubini step has two fixed compact supports and integrable time
  kernel. Both resulting operators are bounded on $E$: the left by
  the mean theorem, the right by the bounded smooth $h_{\Psi,s}$
  and compact smooth mask. Dense states extend the identity to $E$.
  For this fixed $\Psi$ the temporal integral exists at every $s$,
  with norm at most $K_0\sup_t\|\Psi_t\|_{H_{\mathrm{sp}}}$.
  Applying $J_R$ gives the explicit version in
  \eqref{E:positive-raw-mean} for every $R,s$.
  For each $R$ it agrees with the completed $L^r$ mean class.
  This makes no simultaneous-version assertion over all rough controls.
  This proves Lemma~\ref{L:positive-raw-embedding}.
\end{proof}

\begin{lemma}[Actual finite labels and entrywise order]
  \label{L:positive-physical-gram}
  Take a finite partition of $I$ with nonempty cells $I_i$,
  lengths $\ell_i>0$ and $\sum_i\ell_i=1$, and real
  $f_i\in C_c^\infty(\R^d)$ with $\|f_i\|_2=1$.
  Put
  \[
    h_i=\mathbf1_{I_i}/\ell_i,\quad b_i=\|\nabla f_i\|_2^2,\quad
    k_i^R=h_i\otimes(m_Rf_i^2),\quad
    k_i^\infty=h_i\otimes f_i^2 .
  \]
  These are elements of $H_I$. If $v_L$ is the multiplier form
  of~\eqref{E:uniform-E-pairing}, then for every $\Phi\in H_I$,
  \begin{equation}\label{E:positive-label-row}
    \langle\Phi,k_i^R\rangle_{H_I}
    =\ell_i^{-1}\int_{I_i}v_{\mathcal T_R\Phi(s)}(f_i,f_i)\ud s .
  \end{equation}
  The Grams $C^R_{ij}=\langle k_i^R,k_j^R\rangle$ and
  $D_{ij}=\langle k_i^\infty,k_j^\infty\rangle$ satisfy
  $0\le C^R_{ij}\le D_{ij}$ entrywise.
\end{lemma}

\begin{proof}
  Lemma~\ref{L:positive-raw-embedding} puts $h_i$ in $H_t$.
  The compact smooth spatial factors lie in $H_{\mathrm{sp}}$
  by its test definition and~\eqref{E:positive-local-spatial-bound}.
  On a smooth finite tensor, use~\eqref{E:uniform-mean-row}
  with $z=(f_i,0)$ and $w=(0,\Lambda f_i)$.
  Its row is $(\eta_Rf_i)(\eta_R\Lambda^{-1}\Lambda f_i)=m_Rf_i^2$.
  Temporal integration gives the pairing with $h_i$.
  This proves~\eqref{E:positive-label-row} on dense tensors.
  Its right side is a bounded raw functional, with bound
  $\ell_i^{-1/r}C_{\rm mean}\|\Phi\|\|f_i\|_{H^{1/2}}^2$
  by Hölder and~\eqref{E:uniform-completed-mean}.
  The left side is bounded by Hilbert Cauchy--Schwarz.
  Thus the identity extends to the completed $H_I$.
  The form is the actual displacement form by
  Lemma~\ref{L:uniform-E-lift}.

  Each Gram is symmetric positive semidefinite.
  Its entries factor into the temporal pairing of nonnegative
  $h_i,h_j$ and the physical spatial covariance pairing.
  The kernels $K,\Gamma$ and $f_i^2,f_j^2$ are nonnegative;
  $0\le m_R\le1$. Tonelli therefore gives the entrywise order.
  The physical pairing already contains $\Gamma$, and the row
  cancellation introduces no additional Fourier factor.
  Zero labels are permitted.
  All hypotheses of Proposition~\ref{P:gram-entrywise-comparison}
  are met, without assuming $D-C^R$ is positive semidefinite.
  This proves Lemma~\ref{L:positive-physical-gram}.
\end{proof}

\begin{lemma}[Uniform optimizer norm and positive smoothing]
  \label{L:positive-uniform-optimizer}
  Put $M=4C_{\rm mean}+4$.
  The supremum defining $\mathcal S_R$ may be restricted to
  $\|\Phi\|\le M$, uniformly in $R$.
  For every finite trial family in Lemma~\ref{L:positive-physical-gram},
  the unmasked finite score has a maximizer $\Phi_D$ with
  $\|\Phi_D\|\le M$. It is zero or is a positive conic combination
  of the $k_i^\infty$. In the latter case it is approximable in $H_I$
  by tests in $\mathcal C_{\mathrm{pos}}$.
\end{lemma}

\begin{proof}
  Put $C=C_{\rm mean}$ and $q=\|\Phi\|$.
  Proposition~\ref{P:uniform-rate-averaging} gives
  \[
    2\operatorname{Rate}_R(\Phi)-q^2/2
    \le2+2Cq-q^2/2<2-2q<0\quad(q>M).
  \]
  The zero control gives zero, proving the restriction.

  Denote the finite unmasked objective by $\mathfrak F_D$.
  Proposition~\ref{P:gram-positive-optimizer} supplies an attained
  maximum at zero or a positive conic sum $\Phi_D$.
  Choose a finite $R'\ge1$ whose plateau contains all supports of $f_i$.
  Then $k_i^{R'}=k_i^\infty$.
  Write $W_i=\ell_i^{-1}\int_{I_i}\mathcal T_{R'}\Phi(s)\ud s$.
  The row identity and real Rayleigh formula imply, for every $\Phi$,
  \begin{align*}
    \mathfrak F_D(\Phi)
    &\le2\sum_i\ell_i\lambda(W_i)-q^2/2\\
    &\le2+2\sum_i\ell_i\|W_i\|-q^2/2
    \le2+2Cq-q^2/2 .
  \end{align*}
  Here $\lambda(L)\le\max(1,\|L\|)\le1+\|L\|$,
  the norm of an average is bounded by the average norm, and
  the $L^1$ mean norm is at most its $L^r$ norm.
  No time-averaging error or concavity of the positive-part square root
  enters this estimate. Since the maximum is nonnegative, its optimizer
  lies in the same ball of radius $M$, independently of the partition,
  profiles and $R'$.

  On a nonzero branch write
  $\Phi_D=\sum_{i\in A}w_i h_i\otimes f_i^2$, with $w_i>0$.
  Extend $h_i$ by zero, convolve with a nonnegative compact smooth
  approximate identity, and multiply by smooth cutoffs
  $0\le\chi_j\le1$ compactly supported in $I$, tending pointwise to one.
  The resulting nonnegative $h_{i,j}\in C_c^\infty(I)$ converge to $h_i$
  in $L^2$: the error is bounded by the convolution error plus
  $\|(\chi_j-1)h_i\|_2$, which tends to zero by domination.
  Set $\Psi_j=\sum_{i\in A}w_i h_{i,j}f_i^2$.
  Finite tensor estimates and~\eqref{E:positive-bochner-inclusion}
  give $\Psi_j\to\Phi_D$ in $H_I$.
  The positive conic sum has positive squared norm, since at least
  one eligible diagonal is positive and all Gram entries are nonnegative.
  Thus $\Psi_j\ne0$ eventually, so $\Psi_j\in\mathcal C_{\mathrm{pos}}$.
  Costs converge.
  The integrated rate continuity proved in
  Proposition~\ref{P:uniform-rate-averaging} gives
  \begin{equation}\label{E:positive-rate-continuity}
    |\operatorname{Rate}_{R'}(\Psi_j)-\operatorname{Rate}_{R'}(\Phi_D)|
    \le\sqrt{3+C\|\Psi_j\|+C\|\Phi_D\|}
    \sqrt{C\|\Psi_j-\Phi_D\|}\longrightarrow0 .
  \end{equation}
  This estimate is uniform in $R'$.
  It asserts no uniform derivative bound on the smooth approximants.
  The zero branch will use $S_2^{\mathrm{pos}}\ge0$ directly.
  This proves Lemma~\ref{L:positive-uniform-optimizer}.
\end{proof}

\begin{proposition}[Signed scores bounded by the positive score]
  \label{P:positive-signed-comparison}
  For every $R\ge1$,
  \begin{equation}\label{E:positive-signed-comparison}
    \mathcal S_R\le S_2^{\mathrm{pos}} .
  \end{equation}
\end{proposition}

\begin{proof}
  Use uniform partitions with $n\ge2$ cells of length $1/n$.
  They are the $B=1,\delta=1/n$ partitions in
  Proposition~\ref{P:uniform-rate-averaging}, with the empty last cell omitted.
  That proposition gives $e_n\to0$, independent of $R$, such that
  \[
    \left|\operatorname{Rate}_R(\Phi)
    -\sum_i\ell_i\lambda\left(\ell_i^{-1}
    \int_{I_i}\mathcal T_R\Phi(s)\ud s\right)\right|\le e_n
    \quad(\|\Phi\|\le M).
  \]
  Fix $R,\Phi,n$, with $\|\Phi\|\le M$, and denote these averages by $W_i$.
  Real compact smooth unit trials are dense in the real $H^1$ unit sphere:
  truncate by cutoffs tending to one, whose gradient errors tend to zero,
  then mollify in $H^1$ and normalize in $L^2$.
  Continuity of the multiplier form in $H^{1/2}$ and hence $H^1$
  proves
  \[
    \lambda(W_i)=
    \sup_{\substack{f\in C_c^\infty(\R^d;\R)\\\|f\|_2=1}}
    \sqrt{(v_{W_i}(f,f)-\|\nabla f\|_2^2)_+}.
  \]
  The real Rayleigh supremum is nonnegative by
  Proposition~\ref{P:uniform-E-spectral}, so the positive part
  does not change its supremum.
  For $\varepsilon>0$ choose one near-maximizing $f_i$ per cell
  with error at most $\varepsilon$ in the square root.
  This is a finite choice, not a measurable continuum optimizer.
  Lemma~\ref{L:positive-physical-gram} gives
  \[
    2\operatorname{Rate}_R(\Phi)-\|\Phi\|^2/2
    \le\mathfrak F_{C^R}(\Phi)+2e_n+2\varepsilon .
  \]
  Optimize this finite objective and use entrywise Gram comparison:
  \[
    \mathfrak F_{C^R}(\Phi)
    \le\mathscr S(C^R)\le\mathscr S(D)=\mathfrak F_D(\Phi_D).
  \]
  This compares optimized values, not the same individual control.
  Lemma~\ref{L:positive-uniform-optimizer} proves
  $\|\Phi_D\|\le M$ before time averaging is used again.
  Choose $R'$ covering the finite supports as in that lemma.
  The label row, Rayleigh bound and the same uniform averaging estimate give
  \begin{equation}\label{E:positive-finite-optimizer-bound}
    \mathfrak F_D(\Phi_D)
    \le2\operatorname{Rate}_{R'}(\Phi_D)-\|\Phi_D\|^2/2+2e_n .
  \end{equation}

  If $\Phi_D=0$, its finite score is zero.
  Otherwise use its positive smooth approximants $\Psi_j$.
  Lemma~\ref{L:positive-raw-embedding} identifies their actual mean
  with $B_E(m_{R'}h_{\Psi_j,s})$ and their raw cost with $v_{\Psi_j}$.
  Since $h_{\Psi_j,s}\ge0$ and $0\le m_{R'}\le1$, real $H^1$ form order gives
  $\operatorname{Rate}_{R'}(\Psi_j)\le I_{\Psi_j}$.
  Thus
  $2\operatorname{Rate}_{R'}(\Psi_j)-\|\Psi_j\|^2/2
  \le S_2^{\mathrm{pos}}$.
  Rate and cost convergence pass this inequality to $\Phi_D$.
  Together with~\eqref{E:positive-finite-optimizer-bound} and the
  separate zero branch this yields
  \[
    2\operatorname{Rate}_R(\Phi)-\|\Phi\|^2/2
    \le S_2^{\mathrm{pos}}+4e_n+2\varepsilon .
  \]
  All finite choices were made after fixing $n,R,\Phi$.
  The right side is independent of those choices.
  Let $\varepsilon\downarrow0$, take the supremum over the raw ball,
  use Lemma~\ref{L:positive-uniform-optimizer}, and then let $n\to\infty$.
  This proves~\eqref{E:positive-signed-comparison}.
  Finiteness of the positive score was already proved from the
  independent old canonical upper in
  Proposition~\ref{P:variational-positive-scores}, not from this conclusion.
  This proves Proposition~\ref{P:positive-signed-comparison}.
\end{proof}

\begin{proposition}[Cutoff removal for the optimized score]
  \label{P:positive-score-cutoff}
  The possibly nonmonotone mask family satisfies
  \[
    \lim_{R\to\infty}\mathcal S_R=S_2^{\mathrm{pos}}.
  \]
\end{proposition}

\begin{proof}
  Fix $\Psi\in\mathcal C_{\mathrm{pos}}$ before taking $R$ or moment limits.
  Use the explicit mean version of Lemma~\ref{L:positive-raw-embedding}.
  Its rate $\lambda_R(s)$ obeys
  $0\le\lambda_R(s)\le\lambda_{\Psi,s}$ by real form order.
  For each real compact smooth unit $f$, $m_R=1$ on its support for
  every sufficiently large $R$. The Rayleigh formula therefore gives
  \[
    \liminf_{R\to\infty}\lambda_R(s)^2
    \ge\int h_{\Psi,s}f^2-\|\nabla f\|_2^2.
  \]
  Compact trial density and~\eqref{E:variational-real-rayleigh}
  identify the supremum of the right side as $\lambda_{\Psi,s}^2$.
  Hence $\lambda_R(s)\to\lambda_{\Psi,s}$ for these versions.
  The global smooth-source bound
  $H_\Psi=\sup_{s,x}h_{\Psi,s}(x)<\infty$ gives
  $\lambda_R(s)\le\sqrt{H_\Psi}$.
  Dominated convergence yields
  $\operatorname{Rate}_R(\Psi)\to I_\Psi$.
  Neither monotonicity of $m_R$ nor an unmasked ground eigenfunction was used.

  The raw representative and cost $\|\Psi\|^2=v_\Psi$ are independent of $R$.
  Thus
  $\liminf_R\mathcal S_R\ge2I_\Psi-v_\Psi/2$ for each fixed test.
  Given $\varepsilon>0$, choose one fixed test with score greater than
  $S_2^{\mathrm{pos}}-\varepsilon$; the finite supremum admits such a test.
  Take the liminf and then $\varepsilon\downarrow0$.
  Proposition~\ref{P:positive-signed-comparison} supplies the reverse limsup.
  This proves Proposition~\ref{P:positive-score-cutoff}.
\end{proof}

\begin{theorem}[Intrinsic original second-moment constant]
  \label{T:intrinsic-second-moment}
  For every original admissible model with data $(1,0)$ and fixed $x$,
  \begin{equation}\label{E:intrinsic-second-moment-limit}
    \lim_{T\to\infty}T^{-(4-\alpha-a)/(3-\alpha)}
    \log\E|u_\Gamma(T,x)|^2
    =C_2(\Gamma,a)=S_2^{\mathrm{pos}}
    =2^{(4-\alpha)/(3-\alpha)}S_1^{\mathrm{pos}}.
  \end{equation}
  Its complete positive-control formula is
  \begin{align}
    C_2(\Gamma,a)
    &=\sup_{\Psi\in\mathcal C_{\mathrm{pos}}}
    \left\{2\int_I
    \sqrt{\sup_{\substack{f\in H^1(\R^d;\R)\\\|f\|_2=1}}
    \left(\int h_{\Psi,s}f^2-\|\nabla f\|_2^2\right)}
    \ud s-\frac{v_\Psi}{2}\right\},
    \label{E:intrinsic-positive-formula}\\
    h_{\Psi,s}(x)
    &=\int_I|s-t|^{-a}(\Gamma*\Psi_t)(x)\ud t,\qquad
    v_\Psi=\int_{I\times\R^d}\Psi(s,x)h_{\Psi,s}(x)\ud x\ud s.
    \notag
  \end{align}
  Here $\mathcal C_{\mathrm{pos}}$ consists of all nonzero nonnegative
  real $C_c^\infty(I\times\R^d)$ tests, including nonseparable ones.
  The covariance retains its actual amplitude, the kinetic term has
  coefficient one, and the Gaussian action is $v_\Psi/2$.
  For nonzero covariance $0<C_2<\infty$.
  For $\Gamma=0$ the constant and moment rate are zero.
\end{theorem}

\begin{proof}
  Assume first $\Gamma\ne0$.
  Proposition~\ref{P:original-signed-point-upper}, followed by
  Proposition~\ref{P:positive-score-cutoff} with $R=\delta^{-1}+2$,
  bounds the original point limsup by $S_2^{\mathrm{pos}}$.
  Proposition~\ref{P:variational-positive-scores} bounds the original
  point liminf below by this same value.
  Both are for the same canonical Stratonovich solution, exponent,
  original covariance and variance-half action.
  Their source-first identifications retain all mixed, unequal, self
  and empty chains and both phase components with initial data $(1,0)$.
  The original spatial stationarity proved in
  Proposition~\ref{P:raw-block-original-lower} extends both bounds to
  every fixed $x$.
  Thus the limit exists and equals $S_2^{\mathrm{pos}}$.
  Its homogeneity, positivity and finiteness are
  Proposition~\ref{P:variational-positive-scores}.
  Formula~\eqref{E:variational-real-rayleigh} and the definition of the
  original raw variance give~\eqref{E:intrinsic-positive-formula}.

  If $\Gamma=0$, $H_I=\{0\}$, all means, variances and rates vanish,
  $\mathcal S_R=S_c^{\mathrm{pos}}=0$, and $u=1$.
  Both formulas are direct zero identities.
  This proves Theorem~\ref{T:intrinsic-second-moment}.
\end{proof}

The label-one score is auxiliary; the theorem does not identify it with
an actual first-moment limit or with Chen's printed constants.
It neither repairs a source energy normalization nor asserts a Jensen
improvement. The exponent formally tends to $(4-\alpha)/(3-\alpha)$
as $a\downarrow0$; no interchange with the static moment limit is claimed.
The explicit bounds and full calibration and whole-argument verification
remain separate steps.
